\section{Hướng phát triển tương lai}

Để hệ thống đáp ứng các yêu cầu mở rộng và tối ưu hóa tính năng trong thực tiễn, một số định hướng phát triển tiếp theo được đề xuất như sau:

\begin{enumerate}[label=\arabic*.]
    \item Nâng cấp Giao thức Bảo mật chuyển tiếp hoàn hảo (Perfect Forward Secrecy - PFS):
    Hệ thống hiện tại sử dụng thuật toán RSA để mã hóa khóa phiên AES, chưa hỗ trợ PFS. Việc nâng cấp lên thuật toán trao đổi khóa Elliptic Curve Diffie-Hellman (ECDHE) sẽ đảm bảo mỗi phiên giao tiếp sử dụng một khóa độc lập. Cơ chế này ngăn chặn nguy cơ dữ liệu quá khứ bị giải mã ngay cả khi khóa bí mật dài hạn của máy chủ bị lộ.

    \item Tích hợp Mã hóa đầu cuối (End-to-End Encryption - E2EE):
    Bên cạnh mã hóa tầng giao vận Client-Server, việc áp dụng Giao thức Signal (Signal Protocol) sẽ hỗ trợ mã hóa E2EE cho các phiên trò chuyện. Máy chủ sẽ hoạt động với vai trò trung chuyển (Blind Relay) đối với các luồng dữ liệu mã hóa, tăng cường mức độ bảo mật cho thông điệp người dùng.

    \item Triển khai Kiến trúc Phân tán (Distributed Architecture):
    Để xử lý lưu lượng mạng vượt ngưỡng giới hạn của một máy chủ vật lý, kiến trúc nguyên khối (Monolithic) có thể được chuyển đổi sang mô hình cụm phân tán (Cluster/Microservices). Sự kết hợp với Redis Pub/Sub hoặc Apache Kafka làm hệ thống trung chuyển thông điệp (Message Broker) giữa các Node sẽ hỗ trợ phân tán tải và nâng cao khả năng mở rộng (Scalability) của hệ thống.

    \item Phát triển Giao diện người dùng (GUI) và Giao thức Đa phương tiện:
    Chuyển đổi Client từ giao diện dòng lệnh (CLI) sang giao diện đồ họa thông qua bộ công cụ \texttt{Qt} (C++) hoặc \texttt{Electron/Tauri}. Dự án cũng có thể tích hợp tiêu chuẩn \texttt{WebRTC} nhằm cung cấp tính năng gọi thoại và gọi video theo mô hình mạng ngang hàng (P2P), giảm tải băng thông cho máy chủ trung tâm.
\end{enumerate}
